\scanpage{60}
\mspage{II. 8}
% heading and lines 1-6
1.10. \uline{\unsure{Démonstration.|Démonstr.}} \deleted{\unsure{Il suffit|En effet} \unsure{de|on voit de} 1.4.} \inserted{1°)} Existence de $B$.

Soit $k$ la \deleted{clôture} \inserted{fermeture} algébrique du corps premier dans $K$, \deleted{de sorte que $K$ est une extension régulière} \inserted{et soit $X$ un schéma \deleted{de type fini} \inserted{intègre} de type fini} \deleted{extension} sur $k$, \inserted{de corps des fractions $K$,} tel que $A_K$ provienne d'un schéma abélien sur $X$. \deleted{Alors $X$ est géom.} On peut supposer $X$ \deleted{géom.} normal (et $=$ régulier), et \deleted{il est} d'après la déf. de $k$, il est géom. intègre (NB. $k$ est parfait, étant une extension finie du corps premier).

\diagram{\begin{tikzcd}[ampersand replacement=\&]
X\arrow[r,dash]\arrow[d,dash] \& X'\arrow[d,dash]\\
k\arrow[r,dash] \& k'
\end{tikzcd}}

% lines 9-15
\deleted{\unsure{L'hyp.}} Alors l'extension $\overline K$ de $K$ définit un pt géom. $\unsure{\xi|\zeta}$ de $X$, et l'hyp. sur $A_K$ signifie aussi que la représentation de $\pi_1(X,\xi)$ dans $T_\ell(A_\xi)$ est ess\textsuperscript{t} \deleted{abélienne} \inserted{résoluble}. En vertu de \unsure{1.3|1.5}, \deleted{quitte à} passant à un rev. étale \inserted{surjectif} $X'$ de $X$ et à une extension finie $k'$ de $k$, on trouve que \inserted{$A_{X'}$ est \unsure{radiciellement} isogène à un $B_{k'}\otimes_{k'}X'$} \deleted{la représentation de $\pi_1(X',\xi)$ sur $T_\ell(A(\xi))$ est \unsure{définie} par l'intermédiaire d'une représentation de $\pi_1(\operatorname{Spec}(k'),\xi)$, \unsure{donc} i.e. $T_\ell(A_{X'})$ est un faisceau $\ell$-adique géom. trivial sur $X'$. En vertu du résultat déjà cité de \deleted{Inventiones} (\quad), il s'ensuit que l'on peut trouver $B$ sur $k'$ \deleted{tel que} et une isogénie radicielle \deleted{\unsure{$B\otimes$}} $B_{X'}\longrightarrow A_{X'}$,} ce qui achève la démonstration de l'hyp. en car. quelconque.

% lines 24-27
En car. $p>0$, comme $\operatorname{Gal}(\overline k'/k')\simeq\widehat{\mathbb{Z}}$ est abélien, on conclut que l'action de $\operatorname{Gal}(\overline K/K)$ sur $T_\ell(A(\overline K))$ est en fait \uline{abélienne}.

\deleted{\unsure{Démonstr.} \unsure{Lorsque la} caractéristique est nulle, \unsure{on en conclut que}}

% bottom marginal additions and next heading
\margin{Cas de l'action supposée ess\textsuperscript{t} abélienne}
% slanted
1.11. \uline{Bonne réduction} \inserted{aux pts de car. $\ne\ell$} Si $\overline k'$ est la clôture sép. de $k'$ dans $\overline K$, il est bien connu que les hom. can.
% archive sheet number: 59
